\providecommand{\additionalOptionsForClass}{}
\documentclass[
  accentcolor=tud1c,	% Color theme for TUD corporate design
  colorbacktitle,		% Titlepage has colored background for title area
  inverttitle,			% Font color of title on titlepage is inverted
  \additionalOptionsForClass
  %german%,
  %twoside
]{tudexercise}
\usepackage{courier}

\parindent0em

% Vertical space of 2ex between paragraphs.
% Allow LaTeX to increase or decrease this by 0.3ex. This enables
% more dynamic typesetting.
\setlength{\parskip}{2ex plus0.3ex minus0.3ex}

\raggedbottom

\usepackage[ngerman]{babel}
\usepackage[utf8]{inputenc}
\usepackage{listings}
\usepackage{booktabs}
\usepackage{amsmath}
\usepackage{algorithm2e}
\usepackage{hyperref}
\usepackage{xspace}
\usepackage{tabularx}
\usepackage{tikz}
\usepackage{cleveref}
\usepackage{numprint}
\usepackage{paralist}
\usepackage{verbatim}
\usepackage{tocloft} % for manipulating the table of contents
\usepackage{graphics}
\usepackage{csquotes} % for \enquote{}
\usepackage{siunitx}

\usetikzlibrary{shapes}
\usetikzlibrary{calc}
\usetikzlibrary{arrows}
\usetikzlibrary{decorations}

\usepackage{pifont}
\newcommand{\cmark}{\ding{51}\xspace}%
\newcommand{\xmark}{\ding{55}\xspace}%

\usepackage{todonotes}
%\usepackage[disable]{todonotes} % Use this line to hide all todos

\definecolor{commentgreen}{RGB}{50,127,50}
\lstloadlanguages{C++,[gnu]make,bash}
\lstset{language=C++}
\lstset{captionpos=b}
\lstset{tabsize=3}
\lstset{breaklines=true}
\lstset{basicstyle=\ttfamily}
\lstset{columns=flexible}
\lstset{keywordstyle=\color{purple}}
\lstset{stringstyle=\color{blue}}
\lstset{commentstyle=\color{commentgreen}}
\lstset{otherkeywords=\#include}
\lstset{showstringspaces=false}
\lstset{keepspaces=true}
\lstset{xleftmargin=1cm}
\lstset{literate=%
	{Ö}{{\"O}}1
	{Ä}{{\"A}}1
	{Ü}{{\"U}}1
	{ß}{{\ss}}2
	{ü}{{\"u}}1
	{ä}{{\"a}}1
	{ö}{{\"o}}1
	{'}{{\textquotesingle}}1
	{~}{$\sim$}1
}

\newcommand{\superscript}[1]{\ensuremath{^{\textrm{#1}}}}
\newcommand{\subscript}[1]{\ensuremath{_{\textrm{#1}}}}

\newcommand{\setHeader}[1]{
\providecommand{\examheadertitle}{TODO: Titel einbinden}
\renewcommand{\examheadertitle}{#1}
\begin{examheader}
    \examheadertitle
\end{examheader}
}

\newcommand{\hints}[1]{
\paragraph*{Hinweise}
	\begin{itemize}
		\setlength{\itemsep}{0pt}
		#1
	\end{itemize}
}

\newcommand{\optional}{\xspace(optional)}
\newcommand{\experimental}{\xspace(experimentell)}

\usepackage{fancybox}
\newcommand{\optionaltextbox}{
	\bigskip
	\begin{center}
		\ovalbox{\parbox{0.98\textwidth}{Die Klausur kann ohne diese Aufgabe bestanden werden. Wir empfehlen aber sie trotzdem zu bearbeiten.}}
	\end{center}
}
\newcommand{\experimentaltextbox}{
	\bigskip
	\begin{center}
		\ovalbox{\parbox{0.98\textwidth}{Diese Aufgabe wurde neu erstellt und kann noch Fehler und Inkonsistenzen enthalten. Falls euch etwas derartiges auffällt sprecht uns bitte darauf an oder stellt es auf GitHub in den Issuetracker unter \url{https://github.com/Echtzeitsysteme/tud-cppp/issues}}}
	\end{center}
}

\newcommand{\filename}[1]{\texttt{#1}\xspace}
\newcommand{\bspw}{bspw.\xspace}
\newcommand{\dasheisst}{d.\,h.\xspace}
\newcommand{\zB}{z.\,B.\xspace}
\newcommand{\unteranderem}{u.\,a.\xspace}

\newcommand{\winIdeaGroupName}[1]{\textsf{#1}\xspace}
\newcommand{\menuPath}[1]{\emph{#1}\xspace}
\newcommand{\menuSep}[0]{\ensuremath{\to}\;}
\newcommand{\shortcut}[1]{\texttt{#1}\xspace}

\newcommand{\RK}[1]{\todo[]{\textbf{RK:} #1}}
\newcommand{\RKi}[1]{\todo[inline]{\textbf{RK:} #1}}

\newcommand{\ExercisePrefix}[1]{$[$#1$]$ \xspace}
\newcommand{\ExercisePrefixBasics}{\ExercisePrefix{G}}
\newcommand{\ExercisePrefixMemory}{\ExercisePrefix{S}}
\newcommand{\ExercisePrefixObjectOrientation}{\ExercisePrefix{O}}
\newcommand{\ExercisePrefixAdvanced}{\ExercisePrefix{F}}
\newcommand{\ExercisePrefixEmbeddedC}{\ExercisePrefix{C}}
\newcommand{\ExercisePrefixElevator}{\ExercisePrefix{A}}
\newcommand{\ExercisePrefixAdditionalInformation}{\ExercisePrefix{Z}}

\newcommand{\cpppSolutionName}[1]{Ein Lösungsvorschlag kann in \href{./exercises/solutions/#1/.}{\texttt{#1}} gefunden werden.}

\newcommand{\cclicense}{
\begin{center}
    \ovalbox{
        \parbox{0.98\textwidth}
        {
            \begin{minipage}{0.15\textwidth}
                \includegraphics[width=1.0\textwidth]{00_common/cc_license.png}
            \end{minipage}
            \hfill
            \begin{minipage}{0.80\textwidth}
                \begin{small}
                Dieses Werk ist unter einer Creative Commons Lizenz vom Typ Namensnennung - Nicht kommerziell - Keine Bearbeitungen 4.0 International zugänglich. 
                Um eine Kopie dieser Lizenz einzusehen, konsultieren Sie \url{http://creativecommons.org/licenses/by-nc-nd/4.0/} oder 
                wenden Sie sich brieflich an Creative Commons, Postfach 1866, Mountain View, California, 94042 USA.
                \end{small}
            \end{minipage}
        }
    }
\end{center}
}

%% Custom commands for including a listing, based on \lstinputlisting

% The cpppInputNoPageBreakListing command prevents that Latex inserts a page break in between of the listing.  
% If the listing does not fit onto the page anymore, it will be put onto the next page as a whole.
% The first paramer is optional and arguments from lstinputlisting can be provided here.
% The second parameter is the path to the listing.
\newcommand{\cpppInputNoPageBreakListing}[2][] { \begin{minipage}{\textwidth} \lstinputlisting[#1]{#2} \end{minipage} }
% Variant for inline-inserting code
\newcommand{\cpppNoPageBreakListing}[2][] { \begin{minipage}{\textwidth} \begin{lstlisting}[#1]#2 \end{lstlisting}\end{minipage} }

% The cpppInputListing command allows a listing to continue on the next page.
% The first paramer is optional and arguments from lstinputlisting can be provided here.
% The second parameter is the path to the listing.
\newcommand{\cpppInputListing}[2][] { \vspace{1.5ex} \setbox0\vbox{ \lstinputlisting[#1]{#2} } \unvbox0 }
\newcommand{\cpppListing}[2][] { \vspace{1.5ex} \setbox0\vbox{ \begin{lstlisting}[#1]#2 \end{lstlisting} } \unvbox0 }


